\section{Formal results and proof architecture}
\label{sec:formal}
The preceding sections provide an accessible map of the results. We now record formal statements and principal proofs, with explicit quantifiers and risk guarantees. These formulas should be evaluated within their declared observation families; an asymptotic law derived for one family is not a generic fairness-auditing guarantee.

\paragraph{Observation primitives.}
A Gaussian delta-oracle call reveals only
$Z_v=P_1(v;R)-P_0(v;R)$.  A raw vertex-pair replay reveals the ordered pair
$(P_0(v;R),P_1(v;R))$.  Raw edge, suite, and legal-group replay reveal all
ordered old/new outputs generated with their declared common tape.  Different
rounds use independent tapes.  These primitives can have the same physical
execution count but different sigma-algebras and are not interchangeable in a
converse.  Write $V_+(H)=\{v:\deg_H(v)>0\}$.


\subsection{The Gaussian allocation converse and primal design}


\begin{theorem}[Static upper bound]
\label{thm:upper}
Let $m=|E(H)|\ge1$.  For $\delta\in(0,1)$, a non-adaptive Gaussian
delta-oracle tester exists using at
most
\begin{equation}\label{eq:formal-results-001}
\frac{8\log(2m/\delta)}{\gamma^2}\Psi(H,\nu)+|V_+(H)|
\end{equation}
Gaussian delta-oracle calls and having error probability at most $\delta$.
\end{theorem}

\begin{proof}
Take an optimal solution $q^\star$ and set
\begin{equation}\label{eq:formal-results-002}
n_v=\left\lceil
\frac{8\log(2m/\delta)}{\gamma^2}q_v^\star
\right\rceil.
\end{equation}
For each edge $(u,v)$, the empirical difference of the two vertex means has
variance at most
\begin{equation}\label{eq:formal-results-003}
\frac{\gamma^2}{8\log(2m/\delta)}.
\end{equation}
Reject if an empirical edge
difference has magnitude at least $\gamma/2$.  Under the null, a Gaussian tail
bound followed by a union bound controls all edges.  Under the alternative, an
edge with true difference at least $\gamma$ can be missed only if its estimation
error is at least $\gamma/2$.  The ceilings contribute at most $|V_+(H)|$
Gaussian delta-oracle calls.
\end{proof}


\begin{theorem}[Adaptive lower bound]
\label{thm:lower}
For $\delta\in(0,1/2)$, every adaptive $\delta$-correct tester with finite
expected stopping time (measured in Gaussian delta-oracle calls) has, on the null instance $\theta\equiv0$,
\begin{equation}\label{eq:formal-results-004}
\E_0[T]
\ge
\frac{2\kl(1-\delta,\delta)}{\gamma^2}\Psi(H,\nu).
\end{equation}
\end{theorem}

\begin{proof}
Let $n_v=\E_0[N_v]$.  Fix an edge $(u,v)$ and consider an alternative changing
only its endpoints to $\theta_u=a$ and $\theta_v=-b$, where $a,b\ge0$ and
$a+b=\gamma$.  The adaptive transcript KL from the null to this alternative is
\begin{equation}\label{eq:formal-results-005}
\frac{n_u a^2}{2\nu_u^2}+\frac{n_v b^2}{2\nu_v^2}.
\end{equation}
The standard adaptive change-of-measure argument for sequential sampling
\cite{Chernoff1959,CohenZhao2015}, followed by data processing through the final
accept/reject decision, lower-bounds this by $\kl(1-\delta,\delta)$.  Minimizing over $a+b=\gamma$ yields
\begin{equation}\label{eq:formal-results-006}
\frac{\nu_u^2}{n_u}+\frac{\nu_v^2}{n_v}
\le
\frac{\gamma^2}{2\kl(1-\delta,\delta)}.
\end{equation}
This holds for every edge.  Scaling $n_v$ by
$\gamma^2/(2\kl(1-\delta,\delta))$ therefore gives a feasible solution to
\Cref{eq:psi}, proving the claim.
\end{proof}


\subsection{Replay versus fresh observation}


For the one-edge family, the old thresholds are $p$ at both endpoints. The null applies no differential patch and the alternative adds $d$ to the first endpoint; in four-way replay a common shift $s$ is present in both new endpoint thresholds. These access-specific constructions supply the reference null and alternative distributions used below.


\begin{theorem}[Replay changes small-patch complexity]
\label{thm:bernoulli-separation}
Fix $p\in(0,1)$, $d\in(0,1-p)$, and $\delta\in(0,1/2)$.  For the one-edge family above:

\begin{enumerate}[label=(\roman*)]
\item Under common-uniform raw-pair replay, there is a test using
\begin{equation}\label{eq:formal-results-007}
n_{\rm pair}
=
\left\lceil
\frac{\log \delta}{\log(1-d)}
\right\rceil
\end{equation}
paired rounds with zero type-I error and type-II error at most $\delta$.
Moreover, every adaptive $\delta$-correct paired-replay tester with finite
expected stopping time satisfies, on the null,
\begin{equation}\label{eq:formal-results-008}
\E_0[T]
\ge
\frac{\kl(1-\delta,\delta)}
{(1-p)\log\!\left(\frac{1-p}{1-p-d}\right)}.
\end{equation}

\item Under independent new-version observation, a test using
\begin{equation}\label{eq:formal-results-009}
n_{\rm ind}
=
\left\lceil
\frac{2\log(1/\delta)}{d^2}
\right\rceil
\end{equation}
single-execution new-version calls at $u$ is $\delta$-correct when the retained old baseline
$p$ is known.  Conversely, every adaptive $\delta$-correct independent-replay
tester with finite expected stopping time satisfies
\begin{equation}\label{eq:formal-results-010}
\E_0[T]
\ge
\frac{\kl(1-\delta,\delta)}
{\kl(p,p+d)}.
\end{equation}
\end{enumerate}

Consequently, for every fixed interior $p$,
\begin{equation}\label{eq:formal-results-011}
Q_{\rm pair}(d,\delta)
=
\Theta\!\left(
\frac{\log(1/\delta)}{d}
\right),
\qquad
Q_{\rm ind}(d,\delta)
=
\Theta\!\left(
\frac{\log(1/\delta)}{d^2}
\right)
\end{equation}
as $d\to0$ with $p$ and $\delta$ fixed, up to constants depending on them.
\end{theorem}

\begin{proof}
For paired replay, run $P_0$ and $P_1$ on $u$ using the same
$U\sim\mathrm{Unif}[0,1]$ and reject upon observing a discordant pair
$(Y_0,Y_1)=(0,1)$.  Under $\mathsf H_0$ discordance is impossible.  Under
$\mathsf H_1$ it occurs exactly when $U\in(p,p+d]$, hence independently with
probability $d$ each round.  The miss probability after $n$ rounds is therefore
$(1-d)^n$, proving the upper bound.

For the paired lower bound, one informative query at $u$ has null law
\begin{equation}\label{eq:formal-results-012}
Q_0(1,1)=p,
\qquad
Q_0(0,0)=1-p,
\end{equation}
and alternative law
\begin{equation}\label{eq:formal-results-013}
Q_1(1,1)=p,
\qquad
Q_1(0,0)=1-p-d,
\qquad
Q_1(0,1)=d.
\end{equation}
Thus
\begin{equation}\label{eq:formal-results-014}
\KL(Q_0\|Q_1)
=
(1-p)\log\!\left(\frac{1-p}{1-p-d}\right).
\end{equation}
Queries to $v$ have identical laws in both hypotheses.  The adaptive transcript
chain rule and data processing through the terminal accept/reject decision give
\begin{equation}\label{eq:formal-results-015}
\E_0[N_u]\KL(Q_0\|Q_1)
\ge
\kl(1-\delta,\delta),
\end{equation}
which yields the stated lower bound.

For independent replay, giving the tester the exact old baseline $p$ can only make
the problem easier.  Sampling the new version at $u$ and rejecting when its
empirical mean is at least $p+d/2$ has each one-sided error bounded by
$\exp(-nd^2/2)$ by Hoeffding's inequality, proving the upper bound.
For the lower bound, the old-version observation has the same distribution under
both hypotheses and contributes no information.  Each informative new-version
sample contributes
\begin{equation}\label{eq:formal-results-016}
\KL(\mathrm{Bern}(p)\|\mathrm{Bern}(p+d))
=
\kl(p,p+d).
\end{equation}
The same transcript change-of-measure argument yields the claimed expected-sample
lower bound.

Finally,
\begin{equation}\label{eq:formal-results-017}
(1-p)\log\!\left(\frac{1-p}{1-p-d}\right)
=
d+O(d^2),
\end{equation}
whereas
\begin{equation}\label{eq:formal-results-018}
\kl(p,p+d)
=
\frac{d^2}{2p(1-p)}+O(d^3).
\end{equation}
The asymptotic separation follows.
\end{proof}


\begin{theorem}[Counterfactual replay changes the small-gap exponent]
\label{thm:counterfactual-replay}
Fix \(p>0\), \(s>0\), \(d>0\), \(p+s+d<1\), and
\(\delta\in(0,1/2)\).  The old thresholds are \((p,p)\); the null new
thresholds are \((p+s,p+s)\), and the alternative new thresholds are
\((p+s+d,p+s)\).  Vertex-paired and four-way rounds reveal their ordered raw
outputs, and tapes are independent across rounds.

\begin{enumerate}[label=(\roman*)]
\item Under four-way counterfactual replay, there is a test using
\begin{equation}\label{eq:formal-results-019}
n_{\rm edge}
=
\left\lceil
\frac{\log\delta}{\log(1-d)}
\right\rceil
\end{equation}
replay rounds with zero type-I error and type-II error at most \(\delta\).
Every adaptive \(\delta\)-correct tester with finite expected stopping time
satisfies, on the null,
\begin{equation}\label{eq:formal-results-020}
\E_0[T]
\ge
\frac{\kl(1-\delta,\delta)}
{r\log\!\left(\frac{r}{r-d}\right)},
\qquad
r=1-p-s.
\end{equation}

\item Under vertex-paired replay, even granting the tester \(p\), \(s\), and the
identity of the endpoint that may receive the differential effect, a
\(\delta\)-correct test exists using
\begin{equation}\label{eq:formal-results-021}
n_{\rm vertex}
=
\left\lceil
\frac{2\log(1/\delta)}{d^2}
\right\rceil
\end{equation}
raw vertex-pair rounds.  Conversely, every adaptive \(\delta\)-correct
vertex-paired tester with finite expected stopping time satisfies
\begin{equation}\label{eq:formal-results-022}
\E_0[T]
\ge
\frac{\kl(1-\delta,\delta)}
{
s\log\!\left(\frac{s}{s+d}\right)
+
r\log\!\left(\frac{r}{r-d}\right)
}.
\end{equation}
\end{enumerate}

Consequently, for fixed \(p,s,r>0\),
\begin{equation}\label{eq:formal-results-023}
Q_{\rm edge}(d,\delta)
=
\Theta\!\left(
\frac{\log(1/\delta)}{d}
\right),
\qquad
Q_{\rm vertex}(d,\delta)
=
\Theta\!\left(
\frac{\log(1/\delta)}{d^2}
\right)
\end{equation}
as \(d\to0\) with \(\delta\) fixed, up to constant factors.  Counting concrete program executions
multiplies the two costs by four and two, respectively, and therefore does not
change the exponent separation.
\end{theorem}

\begin{proof}
Under four-way replay, define the pathwise differential
\begin{equation}\label{eq:formal-results-024}
W
=
\bigl(P_1(u;U)-P_0(u;U)\bigr)
-
\bigl(P_1(v;U)-P_0(v;U)\bigr).
\end{equation}
Under \(\mathsf H_0\), the two endpoint deltas are identical for every \(U\), so
\(W=0\) almost surely.  Under \(\mathsf H_1\),
\begin{equation}\label{eq:formal-results-025}
W=1
\quad\Longleftrightarrow\quad
U\in(p+s,p+s+d],
\end{equation}
which occurs with probability \(d\).  Rejecting upon the first \(W=1\) therefore
has zero type-I error and miss probability \((1-d)^n\), proving the upper bound.

For the lower bound, the four-output observation has the same mass \(p\) on the
all-one state and the same mass \(s\) on the state in which both new-version
outputs are one while both old-version outputs are zero.  Under the null, the
remaining all-zero state has mass \(r=1-p-s\).  Under the alternative, mass \(d\)
is split from that state into a new state witnessing the counterfactual
differential.  Hence the null-to-alternative KL of one four-way round is
\begin{equation}\label{eq:formal-results-026}
K_{\rm edge}
=
r\log\!\left(\frac{r}{r-d}\right).
\end{equation}
Adaptive transcript change of measure followed by data processing through the
terminal decision gives
\begin{equation}\label{eq:formal-results-027}
\E_0[T]K_{\rm edge}
\ge
\kl(1-\delta,\delta).
\end{equation}

For vertex-paired replay, querying the potentially changed endpoint \(u\) produces
the pair distribution
\begin{equation}\label{eq:formal-results-028}
(p,\ s,\ r)
\end{equation}
on states \((1,1),(0,1),(0,0)\) under the null, and
\begin{equation}\label{eq:formal-results-029}
(p,\ s+d,\ r-d)
\end{equation}
under the alternative.  Its KL is therefore
\begin{equation}\label{eq:formal-results-030}
K_{\rm vertex}
=
s\log\!\left(\frac{s}{s+d}\right)
+
r\log\!\left(\frac{r}{r-d}\right).
\end{equation}
Queries to \(v\) carry no information, so the same transcript argument proves the
lower bound.  For the upper bound, ignore the extra pair information and retain
only the discordance indicator.  It is Bernoulli\((s)\) under the null and
Bernoulli\((s+d)\) under the alternative.  Thresholding its empirical mean at
\(s+d/2\) and applying Hoeffding's inequality bounds each one-sided error by
\(\exp(-nd^2/2)\).

Finally,
\begin{equation}\label{eq:formal-results-031}
K_{\rm edge}
=
d+O(d^2),
\end{equation}
while the first-order terms in \(K_{\rm vertex}\) cancel and
\begin{equation}\label{eq:formal-results-032}
K_{\rm vertex}
=
\frac{d^2}{2}
\left(
\frac{1}{s}+\frac{1}{r}
\right)
+
O(d^3).
\end{equation}
The claimed asymptotic rates follow.
\end{proof}


\subsection{Threshold nulls and suite synchronization}

For P18, fix $p>0$, $s\ge0$, $\tau\ge0$, and $d>0$ with
$p+s+\tau+d<1$.  The old thresholds are $(p,p)$; the boundary-null new
thresholds are $(p+s+\tau,p+s)$ and the alternative thresholds are
$(p+s+\tau+d,p+s)$.  Each round reveals the ordered raw four-way transcript.


\begin{theorem}[Thresholded replay phase transition]
\label{thm:threshold-replay-phase}
Let $\delta\in(0,1/2)$ and set $r=1-p-s-\tau$.  One four-way synchronized
replay round has null-to-alternative KL
\begin{equation}\label{eq:formal-results-033}
K_\tau(d)
=
\tau\log\!\left(\frac{\tau}{\tau+d}\right)
+
r\log\!\left(\frac{r}{r-d}\right),
\end{equation}
where the first term is zero when $\tau=0$.  Every adaptive
$\delta$-correct tester with finite null expected stopping time satisfies
\begin{equation}\label{eq:formal-results-034}
\E_0[T]\ge
\frac{\kl(1-\delta,\delta)}{K_\tau(d)}.
\end{equation}
Conversely, the full composite problem
\begin{equation}\label{eq:formal-results-035}
\mathsf H_0:\ |p_{1,u}-p_{1,v}|\le\tau,
\qquad
\mathsf H_1:\ |p_{1,u}-p_{1,v}|\ge\tau+d
\end{equation}
admits a fixed-design synchronized-replay test using at most
\begin{equation}\label{eq:formal-results-036}
\left\lceil
\frac{8(\tau+d)\log(1/\delta)}{d^2}
\right\rceil
\end{equation}
rounds.  If $r-d$ is bounded below by a positive constant, then as $d\to0$
with $\delta$ fixed these bounds give
\begin{equation}\label{eq:formal-results-037}
Q_{\rm thr}(\tau,d,\delta)
=
\Theta\!\left(
\frac{\tau+d}{d^2}\log\frac1\delta
\right)
\end{equation}
up to constants.  Hence $\tau=0$ and, more generally, $\tau=O(d)$ retain
$d^{-1}$ scaling, whereas every fixed $\tau>0$ restores $d^{-2}$ scaling as
$d\to0$.
\end{theorem}

\begin{proof}
Under the boundary pair, only two four-way transcript atoms change probability:
masses $\tau$ and $r$ under the null become $\tau+d$ and $r-d$ under the
alternative, giving the displayed KL.  Adaptive transcript change of measure
and data processing yield the lower bound.  Equivalently, conditioning on the
event $U>p+s$ reduces the informative transcript exactly to a Bernoulli shift
from $\tau/(1-p-s)$ to $(\tau+d)/(1-p-s)$.

For the upper bound, in each independent synchronized round retain only
\begin{equation}\label{eq:formal-results-038}
W=\mathbf 1\{P_1(u;U)\ne P_1(v;U)\}.
\end{equation}
Under common-uniform threshold replay, $W$ is Bernoulli with mean equal to the
new-version endpoint gap.  Reject when its empirical mean is at least
$\tau+d/2$.  Binomial stochastic ordering reduces the two composite errors to
the boundary means $\tau$ and $\tau+d$; Bernstein/Chernoff bounds make both at
most $\delta$ at the stated sample size.  Finally, for fixed positive $\tau$
and $r$,
\begin{equation}\label{eq:formal-results-039}
K_\tau(d)
=
\frac{d^2}{2}\left(\frac1\tau+\frac1r\right)+O(d^3),
\end{equation}
whereas $K_0(d)=d+O(d^2)$.  More uniformly, if $r-d\ge\rho>0$,
$K_\tau(d)=O_\rho(d^2/(\tau+d))$, matching the upper order.
\end{proof}


\begin{theorem}[Synchronized suite replay]
\label{thm:suite-replay}
Let \(H=(V,E)\) contain at least one fairness edge and fix
\(p,s,d>0\) with \(p+s+d<1\) and \(\delta\in(0,1/2)\).  Under the null,
every active vertex has old threshold \(p\) and new threshold \(p+s\); under
an alternative, exactly one endpoint of an active edge has new threshold
\(p+s+d\). A raw suite round
reveals both version outputs at every active vertex under one common tape.
There is a test that distinguishes the null from the union of all
single-changed-endpoint alternatives using
\begin{equation}\label{eq:formal-results-040}
T_{\rm suite}
=
\left\lceil
\frac{\log\delta}{\log(1-d)}
\right\rceil
\end{equation}
synchronized suite-replay rounds, with zero type-I error and type-II error at
most \(\delta\) for every alternative.  Hence its concrete execution cost is
\begin{equation}\label{eq:formal-results-041}
2n_+
\left\lceil
\frac{\log\delta}{\log(1-d)}
\right\rceil.
\end{equation}
Moreover, every adaptive \(\delta\)-correct tester whose primitive observation is
a synchronized suite-replay round and whose stopping time has finite expectation
satisfies, on the null,
\begin{equation}\label{eq:formal-results-042}
\E_0[T]
\ge
\frac{\kl(1-\delta,\delta)}
{r\log\!\left(\frac{r}{r-d}\right)},
\qquad
r=1-p-s.
\end{equation}
Consequently, for fixed interior \(p,s,r\),
\begin{equation}\label{eq:formal-results-043}
Q_{\rm suite}^{\rm rounds}
=
\Theta\!\left(\frac{\log(1/\delta)}{d}\right),
\qquad
Q_{\rm suite}^{\rm exec}
=
\Theta\!\left(
\frac{n_+\log(1/\delta)}{d}
\right)
\end{equation}
as \(d\to0\) with \(\delta\), the graph, and hence \(n_+\) fixed.  The round bound has no explicit \(|E|\) factor; the concrete
execution cost remains $2n_+$ per round and can scale with the graph.
\end{theorem}

\begin{proof}
In one suite round, form the old/new delta at every active vertex and compare
those deltas across every fairness edge.  Under the null, all vertex deltas are
identical pathwise because every vertex uses the same \(U\), old threshold
\(p\), and new threshold \(p+s\).  Therefore the tester never rejects.

Under an alternative, let \(u\) be the endpoint receiving the additional shift
\(d\).  Whenever
\begin{equation}\label{eq:formal-results-044}
U\in(p+s,p+s+d],
\end{equation}
the delta at \(u\) equals one while every unchanged neighboring endpoint has
delta zero.  At least one incident fairness edge therefore witnesses a
differential.  This event has probability \(d\) in every suite round, so the
probability of missing it for \(T\) rounds is \((1-d)^T\).  The stated upper
bound follows.

For the lower bound, fix any one legal alternative.  Under the null, the suite
observation has an all-common state of mass \(r=1-p-s\).  Under the fixed
alternative, mass \(d\) is split from this state into a new state identifying the
differential endpoint; all other masses are unchanged.  Hence the
null-to-alternative KL of one suite round is
\begin{equation}\label{eq:formal-results-045}
K_{\rm suite}
=
r\log\!\left(\frac{r}{r-d}\right).
\end{equation}
The standard adaptive transcript change-of-measure argument and data processing
through the terminal decision give
\begin{equation}\label{eq:formal-results-046}
\E_0[T]K_{\rm suite}
\ge
\kl(1-\delta,\delta).
\end{equation}
Finally \(K_{\rm suite}=d+O(d^2)\), proving the asymptotic statement.  Each suite
round performs two version executions at each of the \(n_+\) active inputs.
\end{proof}


\subsection{Partial replay structural and achievable geometry}


\begin{proposition}[Separation of the P17 LP geometries]
\label{prop:p17-geometry-separation}
The two fractional-cover geometries need not be comparable within any finite
universal factor.  More precisely: (i) if every fairness edge has endpoints
with identical legal-group incidence signatures, then
$C_{\rm vert}=C_{\rm edge}$; (ii) for $H=K_n$ with the legal groups equal to
its edges and unit group costs,
$C_{\rm vert}=n/2$ and $C_{\rm edge}=\binom n2$, so their ratio is $n-1$;
and (iii) there is a legal partial-replay instance with
$C_{\rm vert}<\infty$ but $C_{\rm edge}=\infty$.
\end{proposition}

\begin{proof}
For (i), identical endpoint incidence makes the edge-witness row
$\mathbf 1\{u,v\in S_i\}_i$ identical to either endpoint's active-vertex
row; equality propagates over each connected fairness component, and duplicate
constraints may be removed.  For (ii), pair groups cover only their own edge, so
every edge variable in the witness LP is at least one.  Summing the $n$
active-vertex constraints counts every pair variable twice and gives
$C_{\rm vert}\ge n/2$, while assigning weight $1/(n-1)$ to every pair
attains equality.  For (iii), take the path $0$--$1$--$2$ with legal groups
$\{0,2\}$ and $\{1,2\}$ at unit cost.  The active-vertex LP has value $2$,
but edge $\{0,1\}$ belongs to no legal group, so the witness LP is infeasible.
\end{proof}


\begin{proposition}[Active-vertex raw-contrast schedule]
\label{prop:p17-active-contrast}
Assume every vertex appearing in a legal replay group has old threshold $p$
and, under the null, new threshold $p+s$.  Under each alternative, exactly one
fairness-active endpoint $x$ instead has new threshold $p+s+d$, with
$p,s,d>0$ and $p+s+d<1$.  Assume also that every legal group in this
proposition contains at least two distinct vertices.  Suppose
$C_{\rm vert}<\infty$.  Let $q^\star$
solve the active-vertex LP, set
\begin{equation}\label{eq:formal-results-047}
L=\left\lceil\frac{\log\delta}{\log(1-d)}\right\rceil,
\qquad
n_i=\lceil Lq_i^\star\rceil,
\end{equation}
and query $S_i$ exactly $n_i$ times with independent tapes across queries.
Within a query, define $D_v=P_1(v;U)-P_0(v;U)$ and reject if the vector
$(D_v:v\in S_i)$ is nonconstant.  The test has zero type-I error and
pointwise type-II error at most $\delta$, with declared cost at most
\begin{equation}\label{eq:formal-results-048}
L C_{\rm vert}+\sum_{i:q_i^\star>0}c_i.
\end{equation}
Consequently the corrected information lower bound and this executable upper
construction use the same active-vertex combinatorial geometry in the specified
homogeneous single-changed-endpoint family.
\end{proposition}

\begin{proof}
Under the null, common-uniform replay gives
$D_v=\mathbf 1\{U\in(p,p+s]\}$ for every member of every legal group, so
the within-group delta vector is constant pathwise and rejection is impossible.
Under an alternative changing $x$, any legal group containing $x$ contains at
least one distinct member $y$ because legal groups have size at least two.
The alternative changes only $x$, hence
\begin{equation}\label{eq:formal-results-049}
D_x\ne D_y
\quad\Longleftrightarrow\quad
U\in(p+s,p+s+d],
\end{equation}
which has probability $d$.  For every active $x$, feasibility of $q^\star$
and upward rounding give
\begin{equation}\label{eq:formal-results-050}
\sum_{i:x\in S_i}n_i
\ge L\sum_{i:x\in S_i}q_i^\star
\ge L.
\end{equation}
Independent tapes across group queries therefore make the miss probability for
alternative $x$ at most $(1-d)^L\le\delta$.  Finally,
$\lceil Lq_i^\star\rceil\le Lq_i^\star+1$ on the positive support,
which gives the stated cost bound.
\end{proof}


\subsection{Heterogeneous matched-control information}


Here $C$ records which legal groups provide same-background control observations and $R$ records boundary-only observations, with positive declared group costs $c$. Every legal group has at least two vertices and reveals its ordered raw old/new transcript under one common tape; rounds are independent. Fix $p,s,\tau,d>0$ with $p+s+\tau+d<1$. Under the simple null all old thresholds equal $p$, low new thresholds equal $p+s$, and high new thresholds equal $p+s+\tau$; alternative $x$ raises only high candidate $x$ by $d$. The following statements do not cover unrestricted heterogeneous populations or simultaneous changes.


\begin{theorem}[Matched controls under a heterogeneous fair null]
\label{thm:control-replay}
Fix \(\delta\in(0,1/2)\) and the declared two-class family above.

\begin{enumerate}[label=(\roman*)]
\item Every adaptive \(\delta\)-correct partial-replay tester with finite null
expected declared cost satisfies
\begin{equation}\label{eq:formal-results-051}
\E_0[\mathrm{cost}]
\ge
\min_{n\ge0}
\left\{
c^\top n:
\bigl(K_{\rm ctrl}C+K_{\rm bnd}R\bigr)n
\ge
\kl(1-\delta,\delta)\mathbf 1
\right\}.
\end{equation}

\item Define the fractional control-cover quantity
\begin{equation}\label{eq:formal-results-052}
C_{\rm ctrl}
=
\min_{q\ge0}
\left\{
c^\top q:
Cq\ge\mathbf 1
\right\}.
\end{equation}
If \(C_{\rm ctrl}<\infty\), a fixed-design tester has zero type-I error and
pointwise type-II error at most \(\delta\) using declared cost at most
\begin{equation}\label{eq:formal-results-053}
L\,C_{\rm ctrl}
+
\sum_{i:q_i^\star>0}c_i,
\qquad
L=
\left\lceil
\frac{\log\delta}{\log(1-d)}
\right\rceil.
\end{equation}

\item For fixed \(p,s,\tau>0\) with \(p+s+\tau<1\),
\begin{equation}\label{eq:formal-results-054}
K_{\rm ctrl}(d)=d+O(d^2),
\qquad
K_{\rm bnd}(\tau,d)
=
\frac{d^2}{2}
\left(
\frac1\tau+\frac1r
\right)
+O(d^3).
\end{equation}
Consequently, if \(C_{\rm ctrl}<\infty\), the lower program and upper
construction have matching leading control-cover geometry and
\(\Theta(d^{-1}\log(1/\delta))\) fixed-confidence order, up to rounding and
standard constant factors.  If some observable high candidate has no
control-compatible legal group, then every tester requires
\(\Omega(d^{-2}\log(1/\delta))\) expected declared cost at fixed
\(\tau>0\).  If a candidate is completely unobservable, no finite-cost
\(\delta\)-correct tester exists.

\item The mixed-topology information floor has an exact two-term singular
expansion.  Let
\begin{equation}\label{eq:formal-results-055}
\begin{aligned}
U&=\{x\in A:C_{xi}=0\ \forall i\},\\
b_2&=\tfrac12(\tau^{-1}+r^{-1}),\\
b_3&=\tfrac13(r^{-2}-\tau^{-2}).
\end{aligned}
\end{equation}
and, for each observable \(x\in U\), let
\(\bar c_x=\min\{c_i:x\in S_i\}\).  Write
\begin{equation}\label{eq:formal-results-056}
C_{\rm ctrl}^{A\setminus U}
=
\min_{q\ge0}
\left\{
c^\top q:C_{A\setminus U}q\ge\mathbf 1
\right\},
\end{equation}
with value zero if \(A\setminus U\) is empty, and let
\(a=\kl(1-\delta,\delta)\).  If every candidate is observable, then as
\(d\downarrow0\) the value \(Q_{\rm LB}(d)\) of the lower program in (i)
satisfies
\begin{equation}\label{eq:formal-results-057}
Q_{\rm LB}(d)
=
\frac{a}{b_2}
\left(\sum_{x\in U}\bar c_x\right)d^{-2}
+
a\left(
C_{\rm ctrl}^{A\setminus U}
-\frac{b_3}{b_2^2}\sum_{x\in U}\bar c_x
\right)d^{-1}
+O(1).
\end{equation}
This is an exact asymptotic for the information lower program, not a matching
upper-bound expansion when \(U\ne\varnothing\).

\item Suppose \(U\ne\varnothing\) and every candidate is observable, and let
\(m=|U|\).  There is a fixed-design hybrid tester with type-I error and every
pointwise type-II error at most \(\delta\), and declared cost at most
\begin{equation}\label{eq:formal-results-058}
N_U\sum_{x\in U}\bar c_x
+L C_{\rm ctrl}^{A\setminus U}
+\sum_{i:q_i^\star>0}c_i,
\end{equation}
where
\begin{equation}\label{eq:formal-results-059}
N_U=
\left\lceil
\frac{8(\tau+d)}{d^2}\log\frac{m}{\delta}
\right\rceil,
\qquad
L=\left\lceil\frac{\log\delta}{\log(1-d)}\right\rceil,
\end{equation}
and \(q^\star\) is an optimizer defining
\(C_{\rm ctrl}^{A\setminus U}\).  Consequently, for every fixed finite
topology and fixed \(\tau>0\), the mixed lower and upper bounds match the
\(d^{-2}\) leading order and the boundary-minimum/control-cover cost geometry
up to constants.  This does not match the exact two-term coefficients or close
the dependence on \(m\) within the broader adaptive expected-cost class; the
next part separates those regimes.

\item The multiplicative \(\log m\) term can be removed from the
\emph{expected null cost} by allowing an unbounded sequential stopping time.
Define the reverse boundary divergence
\begin{equation}\label{eq:formal-results-060}
J_{\rm bnd}(\tau,d)
=
(\tau+d)\log\frac{\tau+d}{\tau}
+(r-d)\log\frac{r-d}{r},
\end{equation}
and
\begin{equation}\label{eq:formal-results-061}
\begin{aligned}
H_0&=\frac{\log(1/\delta)+\log(r/(r-d))}{K_{\rm bnd}(\tau,d)},\\
H_1&=\frac{\log(m/\delta)+\log(1+d/\tau)}{J_{\rm bnd}(\tau,d)}.
\end{aligned}
\end{equation}
There is a sequential hybrid tester with type-I error and every pointwise
type-II error at most \(\delta\), whose expected null declared cost is at most
\begin{equation}\label{eq:formal-results-062}
H_0\sum_{x\in U}\bar c_x
+L C_{\rm ctrl}^{A\setminus U}
+\sum_{i:q_i^\star>0}c_i.
\end{equation}
Under singleton alternative \(x\in U\), its expected boundary-block cost is
at most
\begin{equation}\label{eq:formal-results-063}
H_0\sum_{z\in U\setminus\{x\}}\bar c_z+H_1\bar c_x.
\end{equation}
Thus the boundary block has expected null cost of order
\(d^{-2}\log(1/\delta)\sum_{x\in U}\bar c_x\).
There is no multiplicative \(\log m\); the changed channel still has additive
\(\log m\) delay.  More precisely, if \(Q_0^\star(d)\) denotes the optimal
expected null declared cost in this P19 family and
\(a=\kl(1-\delta,\delta)\), then
\begin{equation}\label{eq:formal-results-064}
\begin{aligned}
\frac{a}{b_2}\sum_{x\in U}\bar c_x
&\le\liminf_{d\downarrow0}d^2Q_0^\star(d)\\
&\le\limsup_{d\downarrow0}d^2Q_0^\star(d)\\
&\le\frac{\log(1/\delta)}{b_2}\sum_{x\in U}\bar c_x.
\end{aligned}
\end{equation}
The remaining multiplicative gap depends only on confidence:
\begin{equation}\label{eq:formal-results-065}
\frac{\log(1/\delta)}{\kl(1-\delta,\delta)}.
\end{equation}
It is independent of topology size and equals approximately
\(1.1305\) at \(\delta=0.05\).  This statement concerns
the exact P19 simple null and singleton alternatives, and does not give a
bounded deterministic horizon.

\item A deterministic declared-cost horizon can be restored without returning
to the fixed-allocation \(m\log m\) penalty.  Fix any
\(\eta\in(0,\delta)\), let \(H_0(\eta)\) and \(H_1(\eta)\) denote the preceding
quantities with \(\delta\) replaced by \(\eta\), and define
\begin{equation}\label{eq:formal-results-066}
B_\eta
=
\frac1{\delta-\eta}
\max_{x\in U}
\left\{
H_0(\eta)\sum_{z\in U\setminus\{x\}}\bar c_z
+H_1(\eta)\bar c_x
\right\}.
\end{equation}
Truncating the sequential boundary block immediately before its declared cost
would exceed \(B_\eta\), and then accepting that block, gives type-I error at
most \(\eta\), pointwise boundary type-II error at most \(\delta\), and the
deterministic total hybrid cost cap
\begin{equation}\label{eq:formal-results-067}
B_\eta
+L C_{\rm ctrl}^{A\setminus U}
+\sum_{i:q_i^\star>0}c_i.
\end{equation}
For unit boundary costs, fixed \(\delta\), and \(\eta=\delta/2\), this cap is
\begin{equation}\label{eq:formal-results-068}
O\!\left(\frac{m+\log m}{d^2}\right)
=O\!\left(\frac{m}{d^2}\right).
\end{equation}
Part~(i) gives a matching \(\Omega(m/d^2)\) lower bound for every deterministic
cap.  Thus adaptive allocation removes the multiplicative \(\log m\) penalty
even with a deterministic horizon, although this Markov truncation is not
claimed to have optimal dependence on \(\delta\) or exact constants.  The
finite-sample certificate can be minimized over the single scalar
\(\eta\in(0,\delta)\); this optimization does not alter the error proof.
\end{enumerate}
\end{theorem}

\begin{proof}
Fix alternative \(x\in A\).  If a queried group contains \(x\) and another
high vertex, the null transcript has a top tape interval of mass \(r\).  Raising
only \(x\)'s new threshold splits mass \(d\) from that interval into a new
\(x\)-specific state, giving \(K_{\rm ctrl}\).  If the group contains \(x\)
but no other high vertex, the existing high--low interval has mass \(\tau\);
the alternative changes masses \((\tau,r)\) to \((\tau+d,r-d)\), giving
\(K_{\rm bnd}\).  Groups omitting \(x\) have zero KL.  Adaptive transcript
change of measure gives one information constraint per \(x\), proving (i).

For (ii), let \(q^\star\) solve the control-cover LP and query group \(S_i\)
\(\lceil Lq_i^\star\rceil\) times.  Within each group, compare old--new patch
deltas only among its high-class members.  Under the null those deltas are
identical pathwise, even though high and low classes differ benignly by
\(\tau\).  Under alternative \(x\), every control-compatible exposure containing
\(x\) becomes nonconstant exactly when
\(U\in(p+s+\tau,p+s+\tau+d]\), an event of probability \(d\).
The cover gives at least \(L\) independent opportunities for each \(x\), so the
miss probability is at most \((1-d)^L\le\delta\).

For (iii), expand the two exact KL expressions.  In the lower LP, scaling
\(z_i=d\,n_i/\kl(1-\delta,\delta)\) makes the coefficient of \(C\) converge to
one and the coefficient of \(R\) converge to zero.  A feasible control-cover
solution gives a uniform upper bound on the scaled optimum.  Because every
\(c_i>0\), any sequence of scaled optimizers therefore has bounded
coordinates; every convergent subsequence has a limit satisfying
\(Cz\ge\mathbf 1\).  Consequently the lower limit of the scaled objective is
at least \(C_{\rm ctrl}\), while the control-cover solution gives the matching
upper limit.  If candidate \(x\) has no control-compatible group, its
information constraint contains only \(K_{\rm bnd}=\Theta(d^2)\), giving the
quadratic lower bound.

For (iv), observe first that a group containing an unmatched candidate
\(x\in U\) cannot contain any other high vertex: otherwise that group would be
a matched control for \(x\).  Hence the LP block for \(x\) is disjoint from
every other high-candidate block and has exact value
\(a\bar c_x/K_{\rm bnd}\).  The Taylor expansion
\begin{equation}\label{eq:formal-results-069}
K_{\rm bnd}(\tau,d)=b_2d^2+b_3d^3+O(d^4)
\end{equation}
therefore contributes
\begin{equation}\label{eq:formal-results-070}
\frac{a\bar c_x}{b_2}d^{-2}
-\frac{ab_3\bar c_x}{b_2^2}d^{-1}+O(1).
\end{equation}
For the covered block, set
\(\epsilon=K_{\rm bnd}/K_{\rm ctrl}=O(d)\) and scale
\(z=K_{\rm ctrl}n/a\).  Its value is
\((a/K_{\rm ctrl})F(\epsilon)\), where \(F(0)\) is precisely
\(C_{\rm ctrl}^{A\setminus U}\).  Positive costs bound every coordinate of
an optimizer uniformly near zero.  Thus the extra term
\(\epsilon Rz\) relaxes each covering constraint by only \(O(\epsilon)\),
and rescaling by \(1+O(\epsilon)\) gives
\(F(\epsilon)=F(0)+O(\epsilon)\).  Since
\(K_{\rm ctrl}=d+O(d^2)\), the covered block contributes
\(aC_{\rm ctrl}^{A\setminus U}d^{-1}+O(1)\).  Summing the disjoint blocks
proves the expansion.

For (v), choose for every \(x\in U\) a cheapest containing group
\(S_{i_x}\).  Because \(x\) has no control-compatible group,
\(S_{i_x}\) contains no other high vertex; since legal groups have size at
least two, it contains a low reference \(y_x\).  On fresh common-uniform
tapes \(R\), form
\begin{equation}\label{eq:formal-results-071}
W_x=\mathbf 1\{P_1(x;R)\ne P_1(y_x;R)\}.
\end{equation}
Under the specified two-class null this is \({\rm Bernoulli}(\tau)\), and under
singleton alternative \(x\) it is \({\rm Bernoulli}(\tau+d)\).  Apply the
P18 threshold test to each \(x\in U\) with error budget \(\delta/m\) and
\(N_U\) rounds.  A union bound controls the boundary block's type-I error by
\(\delta\), while its miss probability under any unmatched alternative is at
most \(\delta/m\).  Independently apply the zero-type-I matched-control
schedule from (ii) to the residual cover on \(A\setminus U\), with pointwise
miss probability at most \(\delta\), and reject if either block rejects.
The declared costs add to the displayed bound.  Groups containing unmatched
candidates cannot contain another high vertex, so they do not create a hidden
cross-coverage term.  The construction relies on the exact equal-high P19 null
and singleton alternatives; it is not valid for an arbitrary heterogeneous
composite null or simultaneous multi-candidate changes.

For (vi), on the cheapest boundary group for each \(x\in U\), let \(L_x(t)\)
be the likelihood ratio of alternative \(x\) against the exact P19 null after
\(t\) fresh rounds.  Process these channels in any fixed order, rejecting if
the active likelihood ratio reaches \(m/\delta\), moving to the next channel
if it reaches \(\delta\), and accepting after every channel has crossed its
lower boundary.  Under the null, Ville's inequality and a union bound give
type-I error at most \(m(\delta/m)=\delta\).  Under alternative \(x\),
\(1/L_x\) is a nonnegative martingale, so the probability of crossing the
lower boundary, and hence the miss probability, is at most \(\delta\).
The channels use disjoint groups and fresh tapes.  Consequently these Ville
bounds remain valid conditionally on the stopping history of every earlier
channel in the fixed processing order; a false upper crossing on an earlier
null channel is already a correct rejection under a singleton alternative and
cannot create a miss.

The log-likelihood increments are bounded.  Their null mean is
\(-K_{\rm bnd}\), and a lower-boundary crossing overshoots by at most
\(\log(r/(r-d))\).  Wald's identity therefore gives
\(\E_0T_x\le H_0\).  Under alternative \(x\), the mean increment is
\(J_{\rm bnd}\), and an upper crossing overshoots by at most
\(\log(1+d/\tau)\), giving \(\E_xT_x\le H_1\).  Summing costs and adjoining
the residual zero-type-I control schedule proves the displayed bounds.  The
exit times are integrable: from every interior likelihood state, a fixed
finite run of identical informative increments exits through one boundary
with probability bounded uniformly away from zero, giving a geometric tail
over such runs.  This also supplies the integrability condition used in
Wald's identity.  Under alternative \(x\), charging every unchanged channel
at its full null stopping cost, including channels that would occur after
\(x\), is an upper bound for every fixed processing order.  The
lower endpoint of the asymptotic bracket is part~(i), because every unmatched
channel contributes \(a\bar c_x/K_{\rm bnd}\); multiplying the lower and upper
bounds by \(d^2\), using \(K_{\rm bnd}=b_2d^2+O(d^3)\), and observing that the
residual control cost is only \(O(d^{-1})\) proves the bracket.

For (vii), run the sequential boundary block from (vi) at internal error
\(\eta\).  Under singleton alternative \(x\), its untruncated declared cost
has expectation at most the expression inside braces defining \(B_\eta\).
Markov's inequality therefore bounds the probability of truncating before the
unbounded test decides by at most \(\delta-\eta\).  Truncation followed by
acceptance cannot increase type-I error, while its pointwise type-II error is
at most \(\eta+(\delta-\eta)=\delta\).  Appending the deterministic residual
control schedule gives the displayed total cap.  Indeed, on the coupled
untruncated path, a pre-query truncation can occur only if the eventual
untruncated declared cost exceeds \(B_\eta\); hence Markov is being applied to
the complete untruncated cost, not to an optional prefix.  With unit costs and
\(\eta=\delta/2\), the numerator defining \(B_\eta\) is
\begin{equation}\label{eq:formal-results-072}
\frac{(m-1)\log(2/\delta)+\log(2m/\delta)}{b_2d^2}
+O(d^{-1}),
\end{equation}
which is \(O(m/d^2)\) at fixed \(\delta\).  Conversely, part~(i) gives null
expected boundary cost at least
\(m\kl(1-\delta,\delta)/K_{\rm bnd}=\Omega(m/d^2)\); every deterministic cost
cap dominates this expectation.
\end{proof}


% Audited October 15 addition: a nuisance-invariant special case and scope fence.
% This is not the missing general P19 composite-null theorem.
\subsection{A nuisance-invariant matched witness and a scope obstruction}
\label{sec:unknown-shared-threshold}

Exact \emph{pathwise} agreement can remove a shared threshold nuisance without
estimating it. Equal \emph{mean} patch effects, however, do not suffice. The
coupling and stopped-KL ingredients are classical
\cite{GlassermanYao1992,NitinawaratAtiaVeeravalli2013}; the proposition records
their consequence for the declared ordered four-output software transcript.
It is a patch-neutrality comparison (or a zero-tolerance edge), not a general
positive-tolerance fairness guarantee or a solution to composite controlled
sensing.

\begin{proposition}[Unknown exactly shared background and a matching obstruction]
\label{prop:unknown-shared-control}
Let one four-way round retain the ordered transcript
\(
X=(P_0(x;U),P_1(x;U),P_0(y;U),P_1(y;U))
\), where \(P_j(z;U)=\mathbf 1\{U\le p_{j,z}\}\), the same
\(U\sim\mathrm{Unif}[0,1]\) is used for all four outputs, and tapes are
independent between rounds. Fix \(\delta\in(0,1/2)\) and known \(d>0\).
\begin{enumerate}[label=(\roman*)]
\item Let
\(
\mathcal N_d=\{(p,h):0<p<h<h+d<1\}
\).
For \((p,h)\in\mathcal N_d\), the null thresholds are
\((p,p)\to(h,h)\), whereas the corresponding alternative is
\((p,p)\to(h+d,h)\). The patch-delta discrepancy witness has probability
zero under every null and exactly \(d\) under every corresponding alternative.
Consequently,
\begin{equation}\label{eq:unknown-shared-horizon}
L(d,\delta)=\left\lceil\frac{\log\delta}{\log(1-d)}\right\rceil
\end{equation}
independent rounds have uniform type-I error zero and type-II error at most
\(\delta\), at a cost of \(4L(d,\delta)\) concrete program executions.

Conversely, consider a sequential tester adapted to the independent-round
transcript filtration, terminating with a decision, whose two errors are at
most \(\delta\) uniformly over the paired null/alternative families. For every
fixed \((p,h)\in\mathcal N_d\), if its null expected stopping time is finite,
then
\begin{equation}\label{eq:unknown-shared-lower}
\E_{0,p,h}[T]\ge
\frac{\kl(1-\delta,\delta)}
{(1-h)\log\bigl((1-h)/(1-h-d)\bigr)}.
\end{equation}
Thus, for fixed \(\delta\) and nuisance restricted to any fixed compact subset
of \(\{(p,h):0<p<h<1\}\), the minimax null-expected round complexity is
\(\Theta(d^{-1})\) as \(d\downarrow0\). This is only a null-side converse;
the reverse single-round KL is infinite.

\item Consider instead the fair null with old thresholds \((p,p+\eta)\) and
new thresholds \((p+s,p+s+\eta)\), where
\begin{equation}\label{eq:mean-matching-domain}
0<p<p+\eta<p+s<p+s+\eta<1.
\end{equation}
Both vertices have expected patch effect \(s\), and both versions have edge
disparity \(\eta\le\tau\). Nevertheless, their tape-wise patch deltas disagree
with probability exactly \(2\eta\) per round. Applying the zero-witness
procedure for \(L(d,\delta)\) rounds to this fair null has false-rejection
probability
\begin{equation}\label{eq:mean-matching-failure}
1-(1-2\eta)^{L(d,\delta)},
\end{equation}
which tends to one for fixed \(\eta>0\) as \(d\downarrow0\). Equality of mean
patch effects therefore cannot replace pathwise matching.
\end{enumerate}
\end{proposition}

\begin{proof}
For part (i), in the coordinate order of the proposition, the complete
single-round laws are
\begin{align}
P_{0,p,h}:&\quad
1111:p,\qquad 0101:h-p,\qquad 0000:1-h,
\label{eq:unknown-shared-null-atoms}\\
P_{1,p,h}:&\quad
1111:p,\qquad 0101:h-p,\qquad 0100:d,\qquad
0000:1-h-d.
\label{eq:unknown-shared-alt-atoms}
\end{align}
The new atom \(0100\) is exactly the patch-delta discrepancy event. This gives
the two witness probabilities and the miss probability \((1-d)^L\).
Moreover, \(\operatorname{supp}(P_{0,p,h})\) is a strict subset of
\(\operatorname{supp}(P_{1,p,h})\), and hence
\begin{equation}\label{eq:unknown-shared-kl-directions}
D(P_{0,p,h}\Vert P_{1,p,h})
=(1-h)\log\!\frac{1-h}{1-h-d},
\qquad
D(P_{1,p,h}\Vert P_{0,p,h})=\infty.
\end{equation}
Under the null, the one-round log-likelihood increment is bounded. The stopped
transcript chain rule and data processing through the terminal decision give
\(\E_{0,p,h}[T]D(P_{0,p,h}\Vert P_{1,p,h})\ge
\kl(1-\delta,\delta)\), proving \Cref{eq:unknown-shared-lower}.
Uniformly when \(1-h\) is bounded away from zero,
\(
D(P_{0,p,h}\Vert P_{1,p,h})=d+d^2/(2(1-h))+O(d^3)
\), while \(L(d,\delta)=\Theta(d^{-1})\) at fixed \(\delta\).

For part (ii), the two patch-delta indicators correspond to intervals
\((p,p+s]\) and \((p+\eta,p+s+\eta]\). Under
\Cref{eq:mean-matching-domain}, their symmetric difference consists of two
disjoint intervals of length \(\eta\), so its measure is \(2\eta\).
Independence between rounds proves \Cref{eq:mean-matching-failure} and its
limit. The endpoint mean shifts are both \(s\), while both the old and new
inter-endpoint disparities equal \(\eta\).
\end{proof}


\begin{proposition}[A fixed-allocation familywise penalty]
\label{prop:boundary-familywise-fixed}
Consider \(m\) disjoint, unit-cost P19 boundary channels, and suppose a
nonadaptive design takes deterministic query counts \(n_1,\ldots,n_m\).
If a test has type-I error and every pointwise type-II error at most
\(\delta\), write \(n_{(1)}\le\cdots\le n_{(m)}\).  Then, for every
\(k\in\{1,\ldots,m\}\),
\begin{equation}\label{eq:formal-results-073}
\begin{aligned}
n_{(k)}&\ge\frac{\log(1+D_\delta k)}{\rho_2},\\
D_\delta&=\frac{(1-2\delta)^2}{\delta(1-\delta)},\\
\rho_2&=\log\!\left(1+d^2(\tau^{-1}+r^{-1})\right).
\end{aligned}
\end{equation}
Consequently,
\begin{equation}\label{eq:formal-results-074}
\sum_{x=1}^m n_x
\ge
\frac1{\rho_2}\sum_{k=1}^m\log(1+D_\delta k)
=\Omega\!\left(\frac{m\log m}{d^2}\right)
\end{equation}
for fixed \(\delta,\tau,r\).  Hence the \(\log m\) factor in the fixed
boundary allocation of part~(v) is not merely a union-bound artifact.  The
proposition does not cover adaptive allocations; part~(vi) exploits precisely
that distinction for expected null cost.
\end{proposition}

\begin{proof}
For one boundary round, the null masses \((\tau,r)\) change to
\((\tau+d,r-d)\), so its chi-square divergence is
\(d^2(1/\tau+1/r)\).  After \(n_x\) independent rounds the divergence is
\(\exp(\rho_2n_x)-1\).  Fix the \(k\) least-sampled channels and mix their
singleton alternatives uniformly.  Disjointness makes the mixture
chi-square divergence from the null
\begin{equation}\label{eq:formal-results-075}
\frac1{k^2}\sum_{j=1}^k
\bigl(\exp(\rho_2n_{(j)})-1\bigr)
\le
\frac{\exp(\rho_2n_{(k)})-1}{k}.
\end{equation}
Data processing through the terminal decision lower-bounds the same divergence
by the Bernoulli decision divergence
\(D_\delta\), because rejection has probability at most \(\delta\) under the
null and at least \(1-\delta\) under the mixture.  Rearranging gives the
rankwise bound; summing over \(k\) proves the claim.
\end{proof}
